% SHARD-ID: AQM-31-WEYL-HEISENBERG-RESEARCH-STATUS-RINGS
% SHARD-TITLE: Weyl-Heisenberg Representation Status over Rings and LCA Systems
% SHARD-SUMMARY: Separates the LCA, locally compact ring, finite Frobenius-ring, and finite Clifford/Weil representation literatures.
% SHARD-SUMMARY: Records which theorems support Stone-von Neumann uniqueness and which only support Pauli/stabilizer or Clifford-extension structure.
% SHARD-SUMMARY: Sets the ground-truth boundary for the rational-function field \(k(t)\) construction.
% SHARD-KEYWORDS: Weyl-Heisenberg, commutative rings, Frobenius rings, LCA groups, Stone-von Neumann, Clifford group, Weil representation

\section{Weyl-Heisenberg Representation Status over Rings and LCA Systems}
\label{sec:weyl-heisenberg-research-status-rings}

\subsection{Status and Source Anchors}

This shard records the source status needed before defining a Weyl-Heisenberg
system for \(k(t)\).  The LCA baseline is Prasad's
\path{references/heisenberg_weil/prasad_0912_0574_source/SvN.tex}, lines
95--107 and 164--190.  The modern LCA quantum-kinematical source is
\path{references/heisenberg_weil/beny_crann_lee_park_youn_2204_08162_source/main.tex},
lines 305--328, 411--444, and 461--473.  The locally compact ring source is
\path{references/heisenberg_weil/bekka_2502_00387_source/CCR-GeneralRings-v5.tex},
lines 287--350, 862--881, 997--1038, and 1048--1150.  The finite
Frobenius-ring Pauli sources are
\path{references/heisenberg_weil/gluesing_luerssen_pllaha_1710_09884_source/StabCodesFrob5.tex},
lines 224--249, 286--323, 353--355, 412--423, and 469--504, and
\path{references/heisenberg_weil/sidana_kashyap_2202_00248_source/EAQECCs_over_rings_v15.tex},
lines 103--145 and 180--197.  The finite generalized-Heisenberg source is
\path{references/heisenberg_weil/solomon_2501_00650_source/TotslipVersion5_Aug_25.tex},
lines 315--318, 904--981, 1127--1136, 1637--1649, and 1677--1718.
Finite Clifford/Weil splitting results are registered from Lysenko, Trias,
Hashimoto--Horibe--Hayashi, Korbelar--Tolar, and Galindo in
\path{references/heisenberg_weil/SOURCES.md}.

\subsection{Lamport Dependency Skeleton}

\[
\begin{array}{rcl}
D1&:&\text{An LCA Weyl system uses }F\times\widehat F\text{, not }F\times F
      \text{ unless }F\simeq\widehat F.\\
L1&:&\text{For an LCA group }F,\ W(x,\gamma)=T_xM_\gamma
      \text{ is the canonical irreducible Weyl representation.}\\
D2&:&\text{For a locally compact ring }R,\text{ a character }\lambda
      \text{ defines CCR by }\lambda(a\cdot b).\\
L2&:&\text{Bekka's strong ring Stone-von Neumann theorem requires
      symmetry and }R\simeq\widehat R.\\
D3&:&\text{For a finite Frobenius ring }R,\ X(a),Z(b)\text{ give a
      Pauli/error basis.}\\
L3&:&\text{Commutation is controlled by the character-symplectic pairing.}\\
D4&:&\text{Finite Clifford/Weil splitting is an additional structure,
      not automatic from the Pauli basis.}\\
T1&:&\text{For }K=k(t),\text{ the honest analytic uniqueness theorem
      applies after choosing a full LCA phase space.}
\end{array}
\]

\subsection{LCA Weyl Systems D1--L1}

Let \(F\) be a locally compact abelian group.  Its Pontryagin dual
\(\widehat F\) is the group of continuous characters \(F\to \mathbb T\).  The
canonical Weyl phase space is
\[
  G_F=F\times\widehat F.
\]
For \((x,\gamma),(x',\gamma')\in G_F\), define
\[
  \sigma_{\rm can}((x,\gamma),(x',\gamma'))=\gamma(x').
\]
This is the source convention of Beny--Crann--Lee--Park--Youn, lines
461--465.

The representation space is \(L^2(F)\).  For \(f\in L^2(F)\),
\[
  (T_xf)(y)=f(y-x),\qquad (M_\gamma f)(y)=\gamma(y)f(y).
\]
Define
\[
  W(x,\gamma)=T_xM_\gamma .
\]
Then
\[
\begin{aligned}
  W(x,\gamma)W(x',\gamma')
  &=T_xM_\gamma T_{x'}M_{\gamma'}                       \\
  &=\gamma(x')T_{x+x'}M_{\gamma+\gamma'}                 \\
  &=\sigma_{\rm can}((x,\gamma),(x',\gamma'))W(x+x',
    \gamma+\gamma').
\end{aligned}
\]
Thus \(W\) is a projective representation with multiplier
\(\sigma_{\rm can}\).  The commutator bicharacter is
\[
  \Delta((x,\gamma),(x',\gamma'))
  =\gamma(x')\overline{\gamma'(x)}.
\]

Prasad's Stone-von Neumann-Mackey theorem, lines 181--190, states that the
canonical representation on \(L^2(F)\) is irreducible and that any continuous
unitary representation of the corresponding Heisenberg group with the standard
central character decomposes as copies of it.  This is the clean theorem we
may use when the phase space is \(F\times\widehat F\).

\subsection{Locally Compact Rings D2--L2}

Let \(R\) be a unital second-countable locally compact ring.  Let
\(\lambda\in\widehat R\) be a unitary character of the additive group of \(R\).
For \(d\ge1\), Bekka defines a pair of unitary representations \(U,V\) of
\(R^d\) to satisfy
\[
  U(a)V(b)=\lambda(a\cdot b)V(b)U(a),\qquad
  a\cdot b=\sum_{i=1}^d a_i b_i .
\]
The Schrödinger pair on \(L^2(R^d)\) is
\[
  (U_{\rm Schr}(a)f)(x)=f(x+a),\qquad
  (V_{\rm Schr}(b)f)(x)=\lambda(x\cdot b)f(x).
\]

The strong theorem in Bekka, lines 319--350, assumes
\[
  \lambda(ab)=\lambda(ba)
\]
and that
\[
  \nabla_\lambda:R\longrightarrow\widehat R,\qquad
  a\longmapsto(x\mapsto\lambda(xa))
\]
is a topological group isomorphism.  Under these hypotheses, every separable
CCR pair decomposes into copies of the Schrödinger pair.  The source lists
local fields, finite fields, finite products of finite fields, \(\mathbb Z/N\),
matrix algebras over local fields, and adeles as examples where the relevant
self-duality hypotheses hold in the stated forms.

The warning for \(k(t)\) appears already in the same source: for a countably
infinite ring with the discrete topology, line 1014 records that no character
can satisfy the isomorphism condition.  Therefore a bare discrete
\(K=k(t)\) cannot be treated as a self-dual locally compact ring in this
strong theorem.  It may still support an explicit algebraic representation,
but the LCA uniqueness theorem must be applied to \(K_{\rm disc}\times
\widehat{K_{\rm disc}}\), to a local completion, or to an adelic completion.

\subsection{Finite Frobenius Rings D3--L3}

Let \(R\) be a finite commutative Frobenius ring with generating character
\(\chi\).  On the Hilbert space with basis \(\{v_x:x\in R^n\}\), define
\[
  X(a)v_x=v_{x+a},\qquad Z(b)v_x=\chi(b\cdot x)v_x
\]
for \(a,b,x\in R^n\), where \(b\cdot x=\sum_i b_i x_i\).  For \(P=X(a)Z(b)\) and
\(P'=X(a')Z(b')\), the cited sources compute
\[
  PP'=\chi(b\cdot a')X(a+a')Z(b+b')
\]
and hence
\[
  PP'=\chi(b\cdot a'-b'\cdot a)P'P.
\]
The symplectic inner product used for stabilizer theory is
\[
  \langle(a,b),(a',b')\rangle_s=b\cdot a'-b'\cdot a.
\]
The Pauli group is obtained by adjoining the necessary scalar roots of unity.
In characteristic \(2\), this scalar group is larger than the characteristic
roots; this is exactly the finite-ring version of the even-phase caveat.

Thus finite Frobenius-ring literature strongly supports the Pauli/error-basis
and stabilizer layer: generating characters identify the ring with its
character module, the \(X/Z\) operators are unitary, and commutation is
controlled by the character-symplectic form.  This is not, by itself, a full
Clifford or Weil representation classification over every commutative
Frobenius ring.

\subsection{Finite Generalized Heisenberg and Clifford Status D4}

Solomon's generalized Heisenberg group starts from finite abelian groups
\(A,B,C\) and a bilinear pairing \(\lambda:A\times B\to C\):
\[
  h(a,b,c)h(a',b',c')
  =
  h(a+a',b+b',c+c'+\lambda(a,b')).
\]
The commutator is
\[
  [h(a,b,c),h(a',b',c')]
  =
  h(0,0,\lambda(a,b')-\lambda(a',b)).
\]
If \(A,B\) are finite modules over a commutative ring \(R\) and \(\lambda\)
is \(R\)-balanced, the quotient by the center inherits \(R\)-module
structure.  A character \(p:C\to\mathbb C^\times\) gives the left
Schrödinger representation on functions on \(A\):
\[
  (h(a,b,c)\cdot f)(x)=p(\lambda(x,b)+c)f(x+a).
\]
This is a direct finite algebraic analogue of translations and modulations.

The Clifford/Weil layer is subtler.  Lysenko constructs canonical
representations for finite Heisenberg groups of odd exponent, explicitly
leaving the even-order case open in the cited introduction.  Trias constructs
Weil modules over coefficient rings for finite or local nonarchimedean fields
of characteristic not \(2\).  Hashimoto--Horibe--Hayashi, Korbelar--Tolar,
and Galindo analyze projective Clifford extensions and splitting.  Galindo's
2026 result states that for a finite abelian group \(A\), the Clifford
extension splits iff \(4\nmid |A|\).  This is important status information:
the obstruction is not a defect of notation but a genuine finite-even
phenomenon.

\subsection{Conclusion for \texorpdfstring{\(k(t)\)}{k(t)} T1}

The sourced status is therefore:
\[
\begin{array}{ll}
\text{LCA groups/local fields:} & \text{full Stone-von Neumann uniqueness is
available.}\\
\text{Locally compact rings:} & \text{strong uniqueness requires
}R\simeq\widehat R\text{ via }\lambda.\\
\text{Finite Frobenius rings:} & \text{Pauli/stabilizer theory is robust.}\\
\text{Finite Clifford/Weil:} & \text{splitting and even-order issues are
genuine extra data.}
\end{array}
\]
For \(K=k(t)\), the next shard must therefore distinguish: the algebraic
discrete \(K\)-label Heisenberg group; the full discrete-LCA system
\(K_{\rm disc}\times\widehat{K_{\rm disc}}\); and the local or adelic
completion where \(K\) becomes arithmetic input inside a self-dual LCA
configuration space.
